\documentclass{article}
\usepackage{amsmath, amssymb, amsthm}
\newtheorem{proposition}{Proposition}
\newtheorem{theorem}{Theorem}
\newtheorem{corollary}{Corollary}
\newtheorem{definition}{Definition}
\newtheorem{remark}{Remark}
\usepackage{graphicx}
\usepackage{hyperref}
\usepackage{geometry}
\usepackage{algorithm}
\usepackage{algpseudocode}
\usepackage{xcolor}
\geometry{margin=1in}

\newcommand{\merw}{EigenTree}
\newcommand{\merwft}{EigenTree-FT}
\newcommand{\psiback}{\psi_{\mathrm{backup}}}

\title{\merwft: Fault-Tolerant Hub Failover\\for Decentralized Bandit Coordination}
\author{Salvatore Rastelli, Andrea Perin}
\date{}

\begin{document}

\maketitle

\begin{abstract}
    \merw{} achieves centralized bandit performance through a decentralized spanning tree
    rooted at an emergent hub node. The hub is a single point of failure: if it crashes,
    the entire coordination protocol collapses. We propose \merwft, an extension that
    tolerates permanent hub failure with zero communication overhead during normal
    operation. The key idea is to designate the hub's direct child with the highest
    eigenvector centrality as a passive backup. The backup mirrors the hub's global state at every
    cycle for free, since it already receives the full downlink broadcast. Upon detecting
    hub silence, the backup promotes itself as the new root, re-routes its sibling nodes
    at depth 1, and resumes the bandit protocol from the last known state. Failover
    completes in $O(\mathrm{diam}(G))$ rounds in general and $O(1)$ rounds when the
    backup has direct graph edges to all former hub siblings. The backup's state is at
    most one cycle stale at the moment of promotion. We formalize the protocol, prove
    correctness of the failover procedure, and characterize the conditions under which the
    backup is a principled substitute for the original hub.
\end{abstract}

\section{Introduction}

\merw{} resolves the tension between decentralization and statistical efficiency in
cooperative multi-agent bandits by constructing an emergent coordinator: the hub node
$h^\star$, identified via eigenvector centrality without any agent observing the
global graph. The hub aggregates observations from all agents, runs the bandit algorithm
on global statistics, and broadcasts decisions back down the spanning tree.

This design has a structural weakness: $h^\star$ is a single point of failure. If the hub
crashes permanently, no other node holds the global statistics, no downlink is broadcast,
and all non-hub agents are left waiting indefinitely. Restarting the full \merw{} pipeline
to elect a new hub is infeasible during an ongoing bandit episode: it requires
$\tau_{\mathrm{init}}$ power iteration rounds plus a max-flooding phase, and the new hub
would start from scratch with no accumulated statistics.

We propose a simple extension, \merwft, that resolves this without any overhead during
normal operation. The core observation is that the hub's direct children already receive
every downlink broadcast in the normal protocol. They therefore hold an exact copy of the
hub's global state $(\hat{n}, \hat{s}, P)$ at the end of every cycle -- a form of passive
replication that costs nothing extra. We designate the direct child with the highest
eigenvector centrality score as the backup hub $b^\star$, announced once at tree construction time.
If $h^\star$ goes silent, $b^\star$ promotes itself, assumes the root role, and the
protocol continues with at most one cycle of disruption.

This pattern is well-established in distributed systems. OSPF elects a Backup Designated
Router (BDR) that mirrors every topology update from the Designated Router; upon DR
failure the BDR promotes itself in one hello-interval \cite{rfc2328}. RSTP pre-identifies
an Alternate Port for every non-root switch and activates it in a single
proposal/agreement round \cite{ieee8021w}. The formal correctness of passive replication
is established by Lamport, Malkhi and Zhou \cite{lamport2009vertical}, who show via
reduction to Paxos that a backup mirroring the primary's state is at most one
synchronization round stale and that failover is safe. \merwft{} applies this pattern
to the aggregation-tree setting of \merw.

\section{Background and Setup}

We assume the same setting as \merw: $N$ agents on an undirected connected graph
$G = (V, E)$, a spanning tree $\mathcal{T}$ rooted at hub $h^\star$ constructed via \merw{}
power iteration and max-flooding, with tree depth $D$ and each node $i$ knowing its
parent $\mathrm{par}(i)$, children $\mathrm{ch}(i)$, and depth $d_i$. Each relay cycle
consists of an uplink phase ($D$ steps), a hub aggregation step, and a downlink phase
($D$ steps). At the end of each downlink, every node at depth 1 holds the hub's global
state $(\hat{n}, \hat{s}, P)$.

We add two assumptions not present in \merw.

\begin{definition}[Fail-stop failure]
    A node may fail permanently at an unknown round $t_f$. After $t_f$, it sends no
    messages and processes no packets. All other nodes are assumed to be live. Failure
    is detectable via timeout: if a node expects a message within a known window and
    receives none, it concludes the sender has failed.
\end{definition}

The timeout-based failure detector is necessary. Anagnostou and Hadzilacos
\cite{anagnostou1993} prove that no purely self-stabilizing protocol can distinguish a
permanently crashed node from a very slow live node in an asynchronous network; an
external failure oracle (here, the cycle-length timeout) is required. In the fully
asynchronous model, Chandra and Toueg \cite{chandratoueg1996} show that even
timeout-based detectors are inherently unreliable: they can only suspect a node, not
certify its failure. In our setting this ambiguity is resolved by the synchronous round
structure of the bandit protocol: a missing uplink packet at a round boundary is
conclusive evidence of failure, not delay.

\begin{definition}[2-connectivity]
\label{def:2conn}
    The underlying graph $G$ is \emph{2-connected}: it remains connected after the
    removal of any single node. Equivalently, $G$ has no cut vertices.
\end{definition}

This assumption is essential for rerouting after any node failure. The routing tree
$\mathcal{T}$ is a spanning tree of $G$, not $G$ itself. During normal operation only
tree edges are used; non-tree edges $E \setminus E(\mathcal{T})$ are idle. When a node
$v$ fails, $\mathcal{T}$ loses $v$ and becomes a forest, but $G \setminus \{v\}$ remains
connected by 2-connectivity. Rerouting therefore always uses non-tree edges of $G$ to
re-establish paths that were broken in $\mathcal{T}$. Each node knows its neighbors in
$G$ (not just its parent and children in $\mathcal{T}$), since neighbors in $G$ are
discovered during the power iteration neighbor-exchange phase.

\begin{remark}
    If $G$ is a tree itself (i.e., $\mathcal{T} = G$), then $G$ has no non-tree edges
    and every node is a cut vertex. Any single node failure permanently disconnects some
    part of the network and no rerouting is possible. The 2-connectivity assumption is
    therefore tight: it is the minimal condition that guarantees rerouting always succeeds
    after a single node failure.
\end{remark}

\section{The Backup Hub}

\subsection{Selection}

During max-flooding, the hub $h^\star$ is the global maximum and never receives a flood
packet from any neighbor: the flood propagates toward $h^\star$, not outward from it.
Nodes route uplink packets toward the hub silently, without announcing their routing
decision. The hub therefore does not know $\mathrm{ch}(h^\star)$ until the first uplink
of the bandit phase, when packets from depth-1 nodes arrive directly.

The hub learns its children from the first uplink: every node $j$ at depth 1 sends its
observations directly to $h^\star$, revealing its identity. Along with each uplink packet,
$j$ includes its local $\tilde{\psi}_j$ value (already computed during power iteration
and stored locally). The hub collects all arriving packets, identifies
$\mathrm{ch}(h^\star)$, and selects the backup:
\[
    b^\star = \arg\max_{j \in \mathrm{ch}(h^\star)} \tilde{\psi}_j.
\]
The hub announces $b^\star$ in the \emph{first downlink}, adding a single identifier
field. All agents store $b^\star$'s identity upon receiving it. The first cycle therefore
runs without a backup; hub failure during this single cycle is not covered by \merwft.

\subsection{Analogy with EIGRP DUAL}

The selection of $b^\star$ is directly analogous to the Diffusing Update Algorithm (DUAL)
used in EIGRP \cite{garcia1993eigrp}, the routing protocol deployed in Cisco networks.
We describe DUAL briefly to make the analogy precise.

\paragraph{DUAL setup.} Each router $i$ maintains a routing table entry for every
destination $d$, consisting of:
\begin{itemize}
    \item $\mathrm{FD}_i(d)$: the \emph{feasible distance}, the best known distance from
        $i$ to $d$ at the time the current successor was chosen.
    \item $\mathrm{RD}_j(d)$: the \emph{reported distance} from neighbor $j$ to $d$,
        as advertised by $j$.
    \item $\mathrm{successor}(i, d)$: the next-hop neighbor currently used to reach $d$.
    \item $\mathrm{FS}(i, d)$: the set of \emph{feasible successors}, pre-computed
        backup next-hops.
\end{itemize}

\paragraph{Feasibility condition.} A neighbor $j$ is a feasible successor for destination
$d$ at router $i$ if:
\[
    \mathrm{RD}_j(d) < \mathrm{FD}_i(d).
\]
This condition guarantees that $j$ is strictly closer to $d$ than $i$ was when $i$ last
updated its route, which is sufficient to ensure loop-freedom: $j$ cannot be routing
through $i$ to reach $d$.

\paragraph{Failover.} When the current successor to $d$ fails, DUAL checks whether
$\mathrm{FS}(i, d)$ is non-empty. If so, the feasible successor with the lowest
$\mathrm{RD}_j(d)$ is promoted to successor immediately, with no query to the rest of
the network. Convergence is sub-second. If no feasible successor exists, DUAL initiates
a diffusing computation (a network-wide query) to find a new route, which takes $O(D)$
rounds in the worst case.

\paragraph{Analogy with \merwft.} In our setting, the ``destination'' is the hub
$h^\star$ and the ``distance'' metric is replaced by eigenvector centrality $\tilde{\psi}$.
The feasibility condition becomes: $b^\star$ is a valid backup for $h^\star$ if
$\tilde{\psi}_{b^\star}$ is the highest among $h^\star$'s direct children, i.e.,
$b^\star$ is the most structurally central node adjacent to $h^\star$ in the tree. This
is loop-free by construction: $b^\star$ is a child of $h^\star$, so it routes toward
$h^\star$ and cannot already be routing through $h^\star$ to reach itself. The promotion
in step 2 of Section~4.2 is the \merw{} analog of DUAL's feasible-successor activation:
immediate, local, and requiring no network-wide query. When no backup exists (hub has
only one child, or all children have equal $\tilde{\psi}$), the protocol falls back to a
full \merw{} re-run, analogous to DUAL's diffusing computation.

\subsection{Passive State Mirroring}

During normal operation, $b^\star$ receives every downlink broadcast from $h^\star$
exactly as any other depth-1 node does. It stores the received state:
\[
    (\hat{n}^{b^\star}, \hat{s}^{b^\star}, P^{b^\star}) \leftarrow (\hat{n}, \hat{s}, P)
\]
at the end of each cycle. No extra messages are sent; no extra computation is performed.
This is passive replication in the sense of Lamport, Malkhi and Zhou
\cite{lamport2009vertical}: the backup mirrors primary state synchronously and for free,
since it already lies on the downlink path.

\begin{proposition}[Backup state staleness]
\label{prop:staleness}
    At the end of cycle $r$, the backup $b^\star$ holds $(\hat{n}, \hat{s}, P)$ equal to
    the hub's state after the hub's own pull in cycle $r$. If the hub fails during cycle
    $r+1$ at any point before completing its downlink, the backup's state reflects all
    observations up to and including the hub's pull in cycle $r$.
\end{proposition}

\begin{proof}
    In cycle $r$, the hub receives the full uplink, updates $(\hat{n}, \hat{s})$ with all
    agents' observations, increments $P$, performs its own pull, and broadcasts the
    updated state in the downlink. $b^\star$ receives this downlink at step $d_{b^\star} =
    1$ (it is at depth 1), so it holds the hub's state after the hub's pull. If the hub
    fails at any point in cycle $r+1$ before completing the downlink of cycle $r+1$,
    the last downlink $b^\star$ received was from cycle $r$. The backup's state is
    therefore at most one full cycle stale.
\end{proof}

\section{Failover Protocol}

\subsection{Failure Detection}

Each depth-1 node, including $b^\star$, expects a downlink packet from $h^\star$ within
$D$ steps of the end of the uplink phase. If no packet arrives, it declares the hub
failed. Since all depth-1 nodes are synchronized (they sent their uplink packets at the
same wall-clock step), they all detect silence simultaneously after the same timeout.

\subsection{Promotion}

Upon detecting hub failure, $b^\star$ executes the following promotion procedure:

\begin{enumerate}
    \item \textbf{Self-promotion}: $b^\star$ sets $d_{b^\star} \leftarrow 0$,
        $\mathrm{par}(b^\star) \leftarrow \text{undefined}$, and begins acting as the
        hub using its stored state $(\hat{n}^{b^\star}, \hat{s}^{b^\star}, P^{b^\star})$.

    \item \textbf{Sibling re-routing}: $b^\star$ must notify all former siblings
        $\mathrm{ch}(h^\star) \setminus \{b^\star\}$ that it is the new hub. In
        $\mathcal{T}$, $b^\star$ is not directly connected to its siblings: their only
        shared parent was $h^\star$, now failed. However, by 2-connectivity of $G$
        (Definition~\ref{def:2conn}), each sibling $s$ has at least one neighbor in
        $G \setminus \{h^\star\}$ other than its own subtree. $b^\star$ floods the
        promotion message $\langle b^\star, \text{NEW-HUB}, d=1 \rangle$ over $G$-edges:
        it sends to all its $G$-neighbors, which forward to their $G$-neighbors, until
        every former sibling is reached. Since $G \setminus \{h^\star\}$ is connected
        (by 2-connectivity), the flood reaches all siblings in at most $\mathrm{diam}(G
        \setminus \{h^\star\})$ steps. Each sibling $s$ upon receiving the message sets
        $\mathrm{par}(s) \leftarrow b^\star$ and $d_s \leftarrow 1$ if it has a direct
        $G$-edge to $b^\star$, or re-routes via an intermediate node otherwise. Subtrees
        below depth 1 are unaffected.

    \item \textbf{Depth update}: $b^\star$ computes the new tree height. Since $b^\star$
        moves from depth 1 to depth 0, its own subtree's depths decrease by 1. However,
        the former siblings $s \in \mathrm{ch}(h^\star) \setminus \{b^\star\}$ now route
        through $b^\star$, so nodes in their subtrees gain one extra hop. The new height
        is:
        \[
            D' = \max\!\bigl(D - 1,\; \max_{s \in \mathrm{ch}(h^\star) \setminus \{b^\star\}} (d_s^{\max} + 1)\bigr),
        \]
        where $d_s^{\max}$ is the depth of the deepest leaf in $s$'s subtree. In the
        common case where all depth-1 subtrees have the same height, $D' = D$. In the
        worst case (e.g., a star graph where the hub has only one child $b^\star$ with no
        subtree and all other agents are direct children of $h^\star$), the siblings
        become children of $b^\star$ at depth 1, giving $D' = 1 = D$, so no increase
        occurs. $b^\star$ broadcasts $D'$ in the first downlink after promotion so all
        agents synchronize on the new cycle length.

    \item \textbf{Resume}: $b^\star$ immediately begins the next cycle as the new hub,
        accepting uplink packets from its children (now including the re-routed siblings).
\end{enumerate}

\begin{proposition}[Failover correctness]
\label{prop:failover}
    After promotion, the network forms a valid spanning tree $\mathcal{T}'$ rooted at
    $b^\star$ with height $D' \leq D + 1$. Every node has a unique path to $b^\star$,
    and no cycles are introduced.
\end{proposition}

\begin{proof}
    Before promotion, $\mathcal{T}$ is a directed tree rooted at $h^\star$. Remove
    $h^\star$ from $\mathcal{T}$. The result is a forest: one subtree rooted at each
    child of $h^\star$, including $b^\star$. Step 2 adds a directed edge from each
    sibling $s \in \mathrm{ch}(h^\star) \setminus \{b^\star\}$ to $b^\star$, merging
    all subtrees into a single tree rooted at $b^\star$. Since the original edges within
    each subtree are unchanged and each subtree was cycle-free, and the new edges are
    directed from former siblings to $b^\star$ (which has no parent), no cycle can form.
    The height increases by at most 1 (the path from the deepest node in a former sibling
    subtree now passes through $b^\star$ rather than $h^\star$, which is the same number
    of hops). Every node has a unique directed path to $b^\star$ by construction.
\end{proof}

\begin{proposition}[Failover completeness]
\label{prop:completeness}
    The promotion procedure completes in $O(\mathrm{diam}(G))$ rounds and requires
    $O(|E|)$ messages in the worst case. If $b^\star$ is directly connected in $G$ to
    all its former siblings, it completes in $O(1)$ rounds and $O(|\mathrm{ch}(h^\star)|)$
    messages.
\end{proposition}

\begin{proof}
    Step 1 is local. Step 2 floods the promotion message over $G \setminus \{h^\star\}$,
    which is connected by 2-connectivity; the flood reaches all nodes in at most
    $\mathrm{diam}(G \setminus \{h^\star\}) \leq \mathrm{diam}(G)$ rounds, touching at
    most $|E|$ edges. If every former sibling has a direct $G$-edge to $b^\star$, step 2
    completes in 1 round with $|\mathrm{ch}(h^\star)| - 1$ messages. Steps 3 and 4 are
    piggy-backed on step 2 at no extra cost.
\end{proof}

\section{Discussion}

\subsection{Why Depth-1 Nodes and Not the Whole Tree}

A natural alternative is to broadcast $b^\star$'s identity and the new depth signal to
the entire tree after promotion, so all nodes update their routing tables. This is
unnecessary. Nodes at depth $\geq 2$ already route toward their parent, which routes
toward $b^\star$ via the unchanged sub-tree structure. The only nodes whose parent
pointer changes are the hub's direct siblings, which are exactly the nodes contacted in
step 2. The rest of the tree is oblivious to the hub failure.

\subsection{Recursive Backup: Tolerating Sequential Hub Failures}

The single-backup limitation dissolves if the promotion procedure is applied
\emph{recursively}. After $b^\star$ promotes itself to hub, it is now a hub with no
pre-designated successor. Rather than pre-computing a chain of backups, the promoted hub
simply re-runs the same selection procedure it observed from $h^\star$: it learns its new
children from the first uplink, selects the highest-$\tilde\psi$ child as the new backup
$b^{\star\star}$, and announces it in the first downlink. The protocol is therefore
self-renewing: every hub, whether original or promoted, establishes its own successor
after one cycle.

\paragraph{Recursive protocol.} Formally, define the backup selection and announcement
procedure as a subroutine $\textsc{SelectBackup}(h)$:
\begin{enumerate}
    \item Hub $h$ collects the first uplink, learning $\mathrm{ch}(h)$ and their
        $\tilde\psi$ values.
    \item $h$ sets $b_h \leftarrow \arg\max_{j \in \mathrm{ch}(h)} \tilde\psi_j$ and
        announces $b_h$ in the first downlink. All agents store $b_h$.
    \item Upon hub failure, $b_h$ runs the failover protocol of Section~4.2, promotes
        itself, and immediately calls $\textsc{SelectBackup}(b_h)$.
\end{enumerate}
This recursion terminates naturally when the current hub has no children, meaning it is
the last surviving node and there is nothing left to coordinate.

\begin{proposition}[Sequential hub failure tolerance]
\label{prop:recursive}
    Under the recursive backup protocol, \merwft{} tolerates any finite sequence of
    hub failures $h^\star = h_0, h_1, h_2, \ldots$ where each $h_{k+1}$ is the backup
    selected by $h_k$. After each failure, exactly one cycle runs without a backup, and
    the protocol resumes with a newly selected successor.
\end{proposition}

\begin{proof}
    By induction on the number of failures. At step 0, $h_0 = h^\star$ selects $b_{h_0}$
    after the first uplink (Proposition~\ref{prop:staleness} applies). If $h_0$ fails,
    $h_1 = b_{h_0}$ promotes via Proposition~\ref{prop:failover} and calls
    $\textsc{SelectBackup}(h_1)$. The first cycle of $h_1$'s tenure runs without a
    backup; thereafter $h_1$ has announced $b_{h_1}$. If $h_1$ fails, $h_2 = b_{h_1}$
    promotes and the argument repeats. At each step the number of live nodes decreases by
    at least one (the failed hub), so the sequence terminates in at most $N-1$ failures.
\end{proof}

\paragraph{Depth under sequential failures.} Each promotion may change the tree height
as analyzed in Section~4.2. After $k$ sequential hub failures the height is
$D^{(k)} \leq D + k$ in the worst case, since each promotion adds at most one level.
In practice the promoted nodes follow the highest-$\tilde\psi$ path downward from
$h^\star$, which on heterogeneous graphs (Barabasi-Albert, Erdos-Renyi) tends to stay
near the dense core of the graph, keeping $D^{(k)}$ close to $D$.

\subsection{Non-Hub Node Failure}

When a non-hub node $v$ at depth $d_v$ fails, the impact is local: only the subtree
rooted at $v$ is disrupted. The hub and all nodes outside this subtree continue operating
normally.

\paragraph{Failure detection.} The correct locus of detection is $v$'s parent
$\mathrm{par}(v)$, not $v$'s children and not the hub. In each round, $\mathrm{par}(v)$
expects a forwarded uplink packet from $v$ within a known window (one hop delay). If no
packet arrives, $\mathrm{par}(v)$ declares $v$ failed. This is strictly faster than
child-driven detection (children would only notice failure when their own uplink finds no
relay, which takes $D - d_v$ hops to manifest) and faster than hub-driven detection
(which requires a full cycle to notice a missing contribution). The parent-driven approach
is supported by Chitnis et al.\ \cite{chitnis2009} and the proactive multicast recovery
of Liao et al.\ \cite{liao2007proactive}, who show that parent monitoring minimizes both
detection latency and false-positive rate compared to child or hub detection.

Once $\mathrm{par}(v)$ detects failure, it notifies $v$'s children
$\mathrm{ch}(v)$ directly via $G$-edges (it knows $\mathrm{ch}(v)$ because $v$'s
children sent uplink packets through $v$ to $\mathrm{par}(v)$ in prior rounds, revealing
their identities). It pushes the re-attachment assignment to each child in one message,
eliminating the need for children to independently scan their neighborhood.

\paragraph{Precomputed swap edges.} Rather than scanning $G$-neighbors at failure time,
each node $c$ precomputes a \emph{swap parent} $u_c$ during tree construction, following
the best-swap-edge framework of Nardelli et al.\ \cite{nardelli2003} and its distributed
variant by Datta et al.\ \cite{datta2013}. The swap parent minimizes the maximum depth
increase across $c$'s entire subtree, not just $c$'s own depth:
\[
    u_c = \arg\min_{u \in \mathcal{N}_G(c) \setminus \{v\}} \max_{\ell \in \mathrm{subtree}(c)} (d_u + 1 + d_\ell - d_c),
\]
where $d_\ell - d_c$ is the depth of $\ell$ within $c$'s subtree. This precomputation
requires one extra message per node during the tree construction phase (each node
broadcasts its depth to its $G$-neighbors), costs $O(|E|)$ total messages, and
eliminates any on-the-fly neighbor scanning at failure time. Re-attachment then takes
exactly 1 round: $\mathrm{par}(v)$ pushes the swap assignments, each child $c$ sets
$\mathrm{par}(c) \leftarrow u_c$ and updates its depth as
\[
    d_c' = d_{u_c} + 1,
\]
and propagates the shift to every descendant $\ell$ in its subtree:
\[
    d_\ell' = d_c' + (d_\ell - d_c).
\]
The whole subtree of $c$ shifts uniformly by the offset $d_{u_c} + 1 - d_c$, so the
relative structure within the subtree is preserved.

\paragraph{Depth increase.} If $v$ was the only node at depth $d_v$ on the path between
some leaf $\ell$ and the hub, re-attachment may increase the depth of $\ell$'s subtree.
In the worst case, if $c$ has no neighbor at depth $d_v - 1$ other than $v$, it must
attach to a node at depth $d_v$, making its own depth $d_v + 1$ and pushing all
descendants down by one level. The new tree depth is
\[
    D' = \max\bigl(D_{\mathrm{rest}},\ \max_{c \in \mathrm{ch}(v)} (d_{u_c} + 1 + d_c^{\max} - d_c)\bigr),
\]
where $D_{\mathrm{rest}}$ is the depth of the unaffected part of the tree and
$d_c^{\max}$ is the depth of the deepest leaf in $c$'s subtree before re-attachment.
The hub learns $D'$ via the next uplink: every node piggybacks its current depth, so
the hub takes the max over all received values. In the extreme case of a chain graph
where $v$ is the unique intermediate node at depth $d_v$, this gives $D' = D + 1$.

\paragraph{Consequence for the protocol.} A depth increase changes the cycle length from
$2D + 1$ to $2D' + 1$. The hub must broadcast the updated $D'$ in the next downlink so
all agents synchronize on the new cycle length. This requires one extra downlink field,
identical to the depth signal already present in the protocol. No other structural change
is needed: the bandit algorithm running at the hub is unaffected, and all live agents
continue routing toward the hub via their updated parent pointers.

\begin{proposition}[Non-hub failover correctness]
\label{prop:nonhub}
    After re-attachment of all orphaned children of a failed node $v$, the remaining live
    nodes form a valid spanning tree rooted at $h^\star$ (or $b^\star$ if hub failover
    has already occurred) with height $D' \leq D + |\mathrm{ch}(v)|$. No cycles are
    introduced.
\end{proposition}

\begin{proof}
    Before failure, $\mathcal{T}$ is a directed tree. Remove $v$: the result is a forest
    with one subtree rooted at each $c \in \mathrm{ch}(v)$ and one subtree containing the
    hub. By 2-connectivity of $G$ (Definition~\ref{def:2conn}), $G \setminus \{v\}$ is
    connected, so each orphaned child $c$ has at least one live $G$-neighbor $u_c \in
    \mathcal{N}_G(c) \setminus \{v\}$ (which may be a tree edge to a sibling subtree or
    a non-tree edge of $G$) with finite depth. Since $G \setminus \{v\}$ is connected and
    the hub is in $G \setminus \{v\}$, $u_c$ lies in the hub's component. We select $u_c$
    with $d_{u_c} < d_c$, which is possible because $u_c$ is not itself an orphan (only
    children of $v$ are orphaned under single node failure) and thus has a well-defined
    depth. Adding the directed edge $c \to u_c$ connects $c$'s subtree to the hub's
    component. No cycle forms: $u_c$ routes toward the hub and not through $c$. Repeating
    for each $c \in \mathrm{ch}(v)$ merges all orphaned subtrees into the hub's component.
    The depth of any node in $c$'s
    subtree increases by at most 1 per re-attachment; since there are $|\mathrm{ch}(v)|$
    orphaned children, $D' \leq D + |\mathrm{ch}(v)|$.
\end{proof}

\begin{remark}
    In practice $|\mathrm{ch}(v)|$ is small (the tree degree of an intermediate node),
    so the depth increase is bounded and modest. On graphs with high connectivity,
    orphaned children typically find a neighbor at the same or lower depth, giving
    $D' = D$ with no cycle length change at all.
\end{remark}

\subsection{Remaining Limitations}

\paragraph{Simultaneous failures.}
If the current hub and its designated backup fail simultaneously, or if the backup fails
before the hub, the protocol has no failover mechanism. Under the recursive protocol,
the unprotected window is exactly one cycle (between hub promotion and backup
announcement); a failure of the new hub during this window leaves the network without a
coordinator. This is a fundamental constraint of the fail-stop model with a single
failure detector: distinguishing simultaneous failure from sequential failure requires
either redundant failure detectors or a Byzantine fault model, both of which are out of
scope.

\paragraph{State staleness and regret.}
Proposition~\ref{prop:staleness} shows the backup's state is at most one cycle stale.
The consequence for bandit performance -- how much regret is incurred by resuming from a
one-cycle-old state -- is left for future work. Informally, UCB's confidence intervals
are monotonically widening in $P$, so using a slightly smaller $P$ makes the algorithm
more exploratory rather than less, which is conservative rather than catastrophic.

\subsection{Relation to the Literature}

The passive backup pattern of \merwft{} is an instance of primary-backup replication,
formalized by Lamport, Malkhi and Zhou \cite{lamport2009vertical} as a special case of
Paxos. The one-cycle staleness bound (Proposition~\ref{prop:staleness}) corresponds
exactly to the one-synchronization-round staleness of passive replication.

The impossibility result of Anagnostou and Hadzilacos \cite{anagnostou1993} establishes
why the timeout-based failure detector in Section~4.1 is necessary and not an
engineering convenience: distinguishing a crashed hub from a slow hub is impossible
without it in an asynchronous model.

The sibling re-rooting procedure in Proposition~\ref{prop:failover} is a special case
of the tree reconfiguration algorithm of Arora \cite{arora1994}, which handles arbitrary
node failure including root failure in $O(n)$ time. \merwft{} restricts the structural
change to depth-1 nodes, but the sibling notification must flood over $G$-edges since
$b^\star$ and its former siblings share no direct tree edges after $h^\star$ is removed.
The flood completes in $O(\mathrm{diam}(G))$ in general and $O(1)$ when $b^\star$ has
direct $G$-edges to all former siblings -- a condition that holds with high probability on
dense graphs such as Barabasi-Albert networks, where the hub and its neighbors form a
dense clique-like core. The precomputed swap-edge framework of Nardelli et al.\
\cite{nardelli2003} and its distributed linear-time variant \cite{datta2013} provide the
right foundation for the non-hub re-attachment rule, replacing on-the-fly neighbor
scanning with a one-round precomputed assignment.

\bibliographystyle{plain}
\begin{thebibliography}{9}

\bibitem{lamport2009vertical}
L.~Lamport, D.~Malkhi, and L.~Zhou.
\newblock Vertical Paxos and primary-backup replication.
\newblock In \emph{Proceedings of PODC}, pp.~27--36, 2009.

\bibitem{anagnostou1993}
E.~Anagnostou and V.~Hadzilacos.
\newblock Tolerating transient and permanent failures.
\newblock In \emph{Proceedings of WDAG}, LNCS 725, pp.~174--188, 1993.

\bibitem{arora1994}
A.~Arora.
\newblock Efficient reconfiguration of trees.
\newblock In \emph{Proceedings of FTRTFT}, LNCS 863, Springer, 1994.

\bibitem{rfc2328}
J.~Moy.
\newblock {OSPF} version 2.
\newblock \emph{RFC 2328}, IETF, 1998.

\bibitem{ieee8021w}
{IEEE}.
\newblock {IEEE} 802.1w: Rapid reconfiguration of spanning tree, 2001.

\bibitem{garcia1993eigrp}
J.J.~Garcia-Luna-Aceves.
\newblock Loop-free routing using diffusing computations.
\newblock \emph{IEEE/ACM Transactions on Networking}, 1(1):130--141, 1993.

\bibitem{nardelli2003}
E.~Nardelli, G.~Proietti, and P.~Widmayer.
\newblock Swapping a failing edge of a single source shortest paths tree is good and fast.
\newblock \emph{Algorithmica}, 35:56--74, 2003.

\bibitem{datta2013}
A.K.~Datta, L.L.~Larmore, L.~Pagli, and G.~Prencipe.
\newblock Linear time distributed swap edge algorithms.
\newblock In \emph{Proceedings of SIROCCO}, LNCS 8179, pp.~162--174, 2013.

\bibitem{chitnis2009}
L.~Chitnis, A.~Dobra, and S.~Ranka.
\newblock Analyzing the techniques that improve fault tolerance of aggregation trees in sensor networks.
\newblock \emph{Journal of Parallel and Distributed Computing}, 69(12):950--960, 2009.

\bibitem{liao2007proactive}
X.~Liao, H.~Jin, Y.~Liu, L.M.~Ni, and D.~Deng.
\newblock A proactive tree recovery mechanism for resilient overlay multicast.
\newblock \emph{IEEE/ACM Transactions on Networking}, 14(6):1729--1741, 2007.

\bibitem{chandratoueg1996}
T.D.~Chandra and S.~Toueg.
\newblock Unreliable failure detectors for reliable distributed systems.
\newblock \emph{Journal of the ACM}, 43(2):225--267, 1996.

\bibitem{burda2009localization}
Z.~Burda, J.~Duda, J.M.~Luck, and B.~Waclaw.
\newblock Localization of the maximal entropy random walk.
\newblock \emph{Physical Review Letters}, 102(16):160602, 2009.

\bibitem{jelasity2007asynchronous}
M.~Jelasity, G.~Canright, and K.~Eng\o-Monsen.
\newblock Asynchronous distributed power iteration with gossip-based normalization.
\newblock In \emph{European Conference on Parallel Processing}, pp.~514--525, 2007.

\end{thebibliography}

\end{document}
